% Auto-generated complete chapter from 04_eu_portugal_germany_mobility_strategy.md
\chapter{04 — EU / Portugal / Germany Mobility Strategy}
\label{complete:root_005}

\begin{sourcemeta}
\textbf{Fuente:} \href{file:///E:/Laboral/04_eu_portugal_germany_mobility_strategy.md}{04\_eu\_portugal\_germany\_mobility\_strategy.md}\\
\textbf{Seccion:} root\\
\textbf{Estado:} canonical\\
\textbf{Rol:} Modulo principal\\
\textbf{Lineas originales:} 694\\
\textbf{Ultima modificacion:} 2026-06-11 12:54\\
\textbf{Deduplicacion conservadora:} 0 bloques repetidos omitidos; 0 caracteres repetidos no reimpresos.
\end{sourcemeta}

\section*{Resumen de orientacion}
This file defines the legal and operational mobility framework for Alexander Woodcock Salomón, a Colombia-based Unity / Real-Time 3D / Technical Visualization candidate.

\section*{Contenido completo depurado}
\addcontentsline{toc}{section}{Contenido completo depurado}




\subsection{1. Purpose and scope}




This file defines the legal and operational mobility framework for Alexander Woodcock Salomón, a Colombia-based Unity / Real-Time 3D / Technical Visualization candidate.




It is not tax, immigration, or legal advice. It is a structured decision framework based on official EU/German sources and reputable tax references available as of \textbf{11 June 2026}.




The goal is to separate:




\begin{itemize}
\item current legal reality;
\item high-probability future EU mobility through Portuguese citizenship;
\item possible German relocation through marriage to a German citizen;
\item recruiter-facing wording that is accurate and defensible;
\item structural trade-offs between Colombia, Portugal, and Germany.
\end{itemize}




\subsection{2. Source-quality audit}




The first Deep Research run for this section did not produce a usable EU mobility file. It returned salary and remote-work content that belonged to section 03. Therefore, this section was rebuilt manually using the verified mobility and tax sources listed below.




\subsubsection{Primary source categories used}




\begin{itemize}
\item EU official guidance on work permits and EU citizen residence.
\item German federal guidance on EU Blue Card, family reunification, freelancing, and health insurance.
\item German residence law references.
\item PwC tax summaries for Germany and Portugal.
\item Portugal tax and social security references from PwC.
\end{itemize}




\subsubsection{Source limitations}




The Portuguese Sephardic citizenship route itself was not re-audited as a legal case file. The user states that the Portuguese passport is highly likely within approximately two years. This document treats that as a \textbf{high-probability future condition}, not as a current right.




The exact Germany family-reunification process depends on the German mission, local immigration office, marital status, documents, and residence location. The analysis therefore treats marriage to a German citizen as a real strategic route, but not as automatic permission to move without paperwork.




\subsection{3. Current source-of-truth status}




\subsubsection{Legal status now}




Alexander is currently based in Colombia and does \textbf{not} currently have EU work authorization.




He may apply internationally as:




\begin{itemize}
\item Colombia-based remote contractor;
\item freelancer / B2B service provider;
\item candidate open to EOR if the employer supports Colombia;
\item future EU-mobile candidate once Portuguese citizenship is completed.
\end{itemize}




He should not state that he is currently authorized to work in the EU.




\subsubsection{High-probability future status}




The user expects a Portuguese passport within approximately two years through the Sephardic Jewish ancestry route. Once a Portuguese passport is actually issued, Alexander becomes an EU citizen with the general right to live and work across EU/EEA member states, subject to local registration, taxation, and social-security rules.




\subsubsection{Germany relationship route}




The user has a German partner. If they legally marry and relocate to Germany, family reunification can become a practical route to residence and work authorization.




The key distinction is:




\begin{itemize}
\item marriage itself is not enough for employment rights abroad;
\item residence status must be processed correctly;
\item once the relevant German residence title is granted for family reunification, employment can generally be permitted.
\end{itemize}




\subsection{4. Mobility-state matrix}




\begin{center}
\scriptsize
\begin{longtable}{>{\raggedright\arraybackslash}p{0.305\linewidth} >{\raggedright\arraybackslash}p{0.305\linewidth} >{\raggedright\arraybackslash}p{0.305\linewidth}}
\toprule
State & Work authorization & Recruiter wording \\
\midrule
\endfirsthead
\toprule
State & Work authorization & Recruiter wording \\
\midrule
\endhead
Current Colombia & No EU work right & Remote contractor from Colombia \\
Portuguese passport pending & No EU work right yet & Future EU citizenship expected \\
Portuguese passport issued & EU work right & EU citizen, authorized in EU \\
Married + Germany residence title & Germany work right & Authorized to work in Germany \\
Employer-sponsored Blue Card & Germany work right & Visa sponsorship required \\
\bottomrule
\end{longtable}
\end{center}




\subsection{5. EU-wide baseline: what changes with Portuguese citizenship}




Official EU guidance states that EU nationals generally do not need a work permit to work anywhere in the EU. It also states that self-employed EU nationals never require a work permit in another EU country.




This means that after the Portuguese passport is issued, Alexander can frame himself as an EU citizen for employment purposes in EU/EEA contexts. That materially improves his access to:




\begin{itemize}
\item Germany;
\item Portugal;
\item Spain;
\item Netherlands;
\item Ireland;
\item Nordic countries;
\item other EU/EEA labor markets.
\end{itemize}




However, EU freedom of movement does not remove local compliance obligations. If Alexander lives in another EU country for more than three months, that country may require residence registration. EU guidance states that after three months in another EU country, registration may be required, although EU citizens can start work without waiting for the registration certificate.




\subsubsection{Practical consequence}




Portuguese citizenship is a major mobility advantage, but it is not equivalent to frictionless global remote employment. It solves the \textbf{right-to-work barrier} inside the EU. It does not solve:




\begin{itemize}
\item income tax residency;
\item social security affiliation;
\item payroll compliance;
\item employer permanent-establishment risk;
\item health-insurance registration;
\item local banking and invoicing setup.
\end{itemize}




\subsection{6. Portugal route}




\subsection{6.1 What Portugal becomes after citizenship}




Once Alexander receives the Portuguese passport, Portugal becomes both:




\begin{itemize}
\item an EU citizenship anchor;
\item a possible low-friction EU base.
\end{itemize}




Portugal is strategically relevant because it can provide:




\begin{itemize}
\item EU work authorization;
\item easier access to EU employers;
\item proximity to Spain and wider European markets;
\item potentially lower cost of living than Germany or Netherlands;
\item Portuguese-language leverage if developed.
\end{itemize}




\subsection{6.2 Rights as a Portuguese citizen}




As a Portuguese citizen, Alexander would have the general right to:




\begin{itemize}
\item live in Portugal;
\item work in Portugal;
\item work in other EU/EEA countries;
\item provide services as self-employed, subject to local rules;
\item move across the EU under free-movement rules.
\end{itemize}




This right begins only when citizenship/passport status is legally completed. Until then, it should not be used as current work authorization.




\subsection{6.3 Living in Portugal while working remotely}




If Alexander becomes resident in Portugal while working for non-Portuguese employers, the issue is no longer mainly immigration. The main issues become:




\begin{itemize}
\item tax residency;
\item Portuguese social security;
\item whether he is an employee or freelancer;
\item whether the foreign employer needs payroll, EOR, or local registration;
\item whether invoices are issued as independent professional services.
\end{itemize}




A foreign employer cannot simply treat Portugal residence as irrelevant if Alexander works there as an employee. If he is a contractor, the setup is simpler, but Portuguese tax and social-security obligations still apply.




\subsection{6.4 Portugal tax residency}




PwC summarizes Portuguese tax residency as follows: an individual is generally considered tax resident if they spend more than 183 days in Portugal during a 12-month period that begins or ends in the relevant fiscal year, or if they maintain a habitual residence in Portugal.




Portugal taxes residents on worldwide income. Non-residents are generally taxed only on Portuguese-source income.




Portugal’s 2026 personal income tax rates are progressive, ranging from 12.5\% to 48\%, with additional solidarity rates above high-income thresholds.




\subsection{6.5 Portugal social security}




For employment, PwC reports standard Portuguese social security contributions of:




\begin{itemize}
\item 11\% employee contribution;
\item 23.75\% employer contribution.
\end{itemize}




For self-employed workers, PwC reports a 21.4\% contribution rate. Relevant remuneration is generally calculated as 70\% of services income or 20\% of product-related income.




Portugal also applies additional contribution exposure to contracting companies when a high percentage of the worker’s income comes from one company. PwC reports 10\% when 80\% or more of the worker’s income comes from one company, and 7\% when 50\% to 80\% comes from one company.




\subsubsection{Practical consequence}




Portugal can work well as a European base, but contractor economics must be calculated carefully. A gross remote salary may look attractive but can lose value through:




\begin{itemize}
\item income tax;
\item social security;
\item accountant costs;
\item currency conversion;
\item invoicing obligations;
\item possible VAT obligations.
\end{itemize}




\subsection{6.6 Portugal registrations likely needed}




For a Portuguese-resident worker or freelancer, likely administrative steps include:




\begin{itemize}
\item NIF / Portuguese tax identification;
\item social security registration;
\item local address or residence record;
\item activity registration if freelancing;
\item invoicing setup;
\item accountant review if billing foreign clients;
\item bank account or EU-compatible payment infrastructure.
\end{itemize}




These steps should be confirmed with a Portuguese accountant once the passport timeline is concrete.




\subsection{6.7 NHR / IFICI caution}




Portugal’s previous Non-Habitual Resident regime was substantially changed. Various post-NHR incentive frameworks have existed or been proposed, but eligibility is technical and cannot be assumed.




For this strategy, Portuguese tax incentives should be treated as \textbf{optional upside}, not as a planning foundation.




Do not build the career strategy on assuming:




\begin{itemize}
\item automatic 20\% flat taxation;
\item automatic NHR eligibility;
\item automatic tax incentives for technical-art or Unity work;
\item automatic benefit from moving to Portugal.
\end{itemize}




Portugal remains valuable primarily because of EU citizenship and market access, not because of guaranteed tax advantages.




\subsection{7. Germany route}




\subsection{7.1 Germany as target market}




Germany is strategically important for Alexander because it combines:




\begin{itemize}
\item strong industrial visualization market;
\item simulation and digital twin demand;
\item automotive and manufacturing ecosystems;
\item XR and enterprise 3D use cases;
\item possible relocation route through a German partner;
\item higher salary ceiling than Portugal or Spain.
\end{itemize}




Germany is not the easiest immediate route because of language, bureaucracy, health insurance, payroll rules, and work authorization. But it may be one of the strongest 12–24 month routes if legal status is solved.




\subsection{7.2 Marriage / family reunification route}




The user’s German-partner context is strategically significant.




If Alexander legally marries his German partner and relocates to Germany through the correct family-reunification process, the route can provide residence and work authorization.




German federal guidance for spouses joining EU citizens states that third-country spouses usually need a family-reunification visa if they are not visa-exempt, and once a corresponding residence title is granted, they are immediately entitled to take up any kind of employment.




\subsubsection{Practical consequence}




Marriage can materially improve the Germany route, but only after:




\begin{itemize}
\item legal marriage is completed;
\item documents are prepared;
\item the correct visa/residence process is followed;
\item German authorities grant the residence title.
\end{itemize}




Until then, Alexander should not claim German work authorization.




\subsection{7.3 What this enables in Germany}




Once residence and work authorization are granted, Germany becomes more realistic for:




\begin{itemize}
\item Unity Developer;
\item Technical Artist;
\item Simulation Developer;
\item Digital Twin Visualization Developer;
\item XR Developer;
\item Tools / Pipeline Developer;
\item technical visualization roles in industrial companies.
\end{itemize}




This can be stronger than applying from Colombia because German employers often prefer or require:




\begin{itemize}
\item local residence;
\item EU or German work authorization;
\item German payroll eligibility;
\item ability to attend hybrid/on-site meetings;
\item health insurance and tax registration.
\end{itemize}




\subsection{7.4 Blue Card route}




The EU Blue Card is an employer-sponsored route for qualified professionals. German federal guidance lists the 2026 general salary threshold as \textbf{€50,700 gross annual salary}. For shortage occupations and young professionals, the 2026 threshold is \textbf{€45,934.20}.




For IT specialists without a formal degree, Germany allows an EU Blue Card route if the person has a specific job offer, meets the lower salary threshold, and has at least three years of relevant IT work experience within the last seven years at a university-graduate level.




\subsubsection{Relevance to Alexander}




The Blue Card is not the primary route right now because:




\begin{itemize}
\item he does not yet have strong formal industry experience;
\item Unity / Technical Artist roles may not always map neatly to shortage occupation categories;
\item employers may prefer candidates already in the EU;
\item German language can matter for non-startup companies.
\end{itemize}




However, the Blue Card becomes relevant if Alexander secures:




\begin{itemize}
\item a German job offer;
\item sufficient salary;
\item clear degree or experience equivalence;
\item role classification accepted by German authorities.
\end{itemize}




\subsection{7.5 Freelance / self-employment route in Germany}




Germany has a freelance residence route for certain self-employed professionals. German federal guidance states that non-EU citizens intending to work freelance in Germany generally need a residence permit for freelance employment.




To qualify, a freelancer must generally show:




\begin{itemize}
\item ability to finance the undertaking;
\item ability to make a living;
\item relevant permissions for regulated professions if applicable;
\item health insurance;
\item business plan and supporting documents;
\item pension provisions if older than 45.
\end{itemize}




The permit is usually initially limited to a maximum of three years. A settlement permit can become possible later if the activity is successful.




\subsubsection{Relevance to Alexander}




This route is theoretically possible but not recommended as the primary path unless he has:




\begin{itemize}
\item paying EU clients;
\item a clear freelance business case;
\item predictable income;
\item German accountant support;
\item strong documentation.
\end{itemize}




For the next 3–6 months, Germany freelance immigration is less practical than remote contracting from Colombia.




\subsection{7.6 Germany health insurance}




Health insurance is compulsory in Germany. German federal guidance distinguishes statutory and private health insurance and notes that proof of health insurance may be required for visa purposes.




For most employees, statutory health insurance applies. Contributions depend on income. Spouses and children may sometimes be included without extra contributions if they are not working.




\subsubsection{Practical consequence}




Germany’s salary ceiling is higher, but net-income modeling must include:




\begin{itemize}
\item health insurance;
\item pension insurance;
\item unemployment insurance;
\item long-term care insurance;
\item income tax;
\item possible church tax;
\item cost of living.
\end{itemize}




\subsection{7.7 Germany tax and social security}




PwC summarizes German tax residence as being triggered by a dwelling available for use or habitual abode. A habitual abode generally occurs when a person is physically present in Germany for more than six months.




German tax residents are generally taxed on worldwide income. Non-residents are generally taxed on German-source income.




PwC summarizes Germany’s 2026 progressive income tax rates as:




\begin{itemize}
\item 0\% up to €12,096 for single taxpayers;
\item 14\%–42\% from €12,097 to €68,429;
\item 42\% up to €277,825;
\item 45\% above €277,826.
\end{itemize}




Germany also has social-security contributions. PwC reports contribution rates for pension, unemployment, health insurance, and long-term care insurance, usually split between employer and employee.




\subsubsection{Practical consequence}




Germany can provide higher gross salaries than Portugal, Spain, or Colombia. But the gross-to-net difference is significant. A German offer should be evaluated by net salary, not headline gross salary.




\subsection{8. EU-wide remote work considerations}




\subsection{8.1 Right to work is not the same as payroll compliance}




Portuguese citizenship or German residence may solve work authorization, but it does not automatically solve remote work compliance.




A worker living in Portugal or Germany while employed by a foreign company may trigger obligations related to:




\begin{itemize}
\item local payroll;
\item income tax withholding;
\item social-security contributions;
\item employer registration;
\item EOR use;
\item permanent establishment risk;
\item local labor law protections.
\end{itemize}




\subsection{8.2 Remote employee vs contractor}




The difference matters.




As an employee, the employer may need to run payroll or use EOR in the country of residence.




As a contractor, the worker generally invoices independently. But if the contractor works like an employee, misclassification risk may arise.




Indicators of employee-like relationship include:




\begin{itemize}
\item fixed schedule;
\item exclusivity;
\item employer-provided tools;
\item daily direct supervision;
\item inability to subcontract;
\item long-term dependency on one client;
\item integration into the company’s internal hierarchy.
\end{itemize}




\subsection{8.3 Permanent establishment risk}




A foreign company may be concerned that a remote employee in Portugal or Germany creates tax exposure for the company. This is especially sensitive if the worker:




\begin{itemize}
\item negotiates contracts;
\item signs deals;
\item manages local business operations;
\item represents the company locally;
\item operates as a senior decision-maker.
\end{itemize}




For Alexander’s likely early roles, permanent-establishment risk is lower if he works as an individual contributor. But employers may still prefer EOR or contractor arrangements for compliance.




\subsection{9. Recruiter-facing work authorization statements}




\subsection{9.1 Current wording — Colombia, no EU passport yet}




Use:




\begin{quote}
I am currently based in Colombia and available for remote contractor/B2B work with international companies. I can also work through an EOR if the company supports Colombia.
\end{quote}




Avoid:




\begin{quote}
I am authorized to work in the EU.
\end{quote}




Avoid:




\begin{quote}
I can work in Germany without sponsorship.
\end{quote}




Avoid:




\begin{quote}
I have Portuguese citizenship.
\end{quote}




\subsection{9.2 Portuguese passport pending}




Use only if relevant:




\begin{quote}
I am in the final/advanced pathway toward Portuguese citizenship and expect EU mobility within approximately two years, but I am currently applying as a Colombia-based remote contractor.
\end{quote}




Avoid overusing this in short applications. It can distract from immediate employability.




\subsection{9.3 Once Portuguese passport is issued}




Use:




\begin{quote}
I am an EU citizen through Portuguese citizenship and authorized to work in EU/EEA countries. I am open to remote, hybrid, or relocation-based roles depending on the position.
\end{quote}




This becomes a strong ATS and recruiter filter advantage for EU roles.




\subsection{9.4 If married and Germany residence title is granted}




Use:




\begin{quote}
I am legally resident in Germany and authorized to work there.
\end{quote}




If still in process, use:




\begin{quote}
I am pursuing family-based relocation to Germany and will require the appropriate residence/work authorization process before starting local employment.
\end{quote}




\subsection{9.5 Blue Card wording}




Use:




\begin{quote}
I would require employer sponsorship for a German EU Blue Card unless another residence route is completed first.
\end{quote}




Do not claim Blue Card eligibility unless the job offer, salary, degree/experience match, and authority requirements are confirmed.




\subsection{10. Colombia vs Portugal vs Germany: structural comparison}




\begin{center}
\scriptsize
\begin{longtable}{>{\raggedright\arraybackslash}p{0.305\linewidth} >{\raggedright\arraybackslash}p{0.305\linewidth} >{\raggedright\arraybackslash}p{0.305\linewidth}}
\toprule
Base & Main advantage & Main risk \\
\midrule
\endfirsthead
\toprule
Base & Main advantage & Main risk \\
\midrule
\endhead
Colombia & Lower cost base & No EU work right \\
Portugal & EU citizenship base & Tax/social setup \\
Germany & Higher salary ceiling & Bureaucracy/language \\
\bottomrule
\end{longtable}
\end{center}




\subsection{10.1 Colombia contractor}




Best for the next 3–6 months.




Strengths:




\begin{itemize}
\item no relocation friction;
\item immediate availability;
\item lower cost of living;
\item attractive to foreign employers seeking LATAM contractors;
\item time-zone compatibility with US markets;
\item can build portfolio while applying.
\end{itemize}




Weaknesses:




\begin{itemize}
\item no EU employee access;
\item lower salary ceiling;
\item weaker legal protection;
\item must manage taxes and social security independently;
\item foreign employers may reject Colombia if they do not support contractor or EOR arrangements.
\end{itemize}




\subsection{10.2 Portugal resident after citizenship}




Best as a 12–24 month EU mobility base.




Strengths:




\begin{itemize}
\item full EU citizenship rights after passport;
\item simpler EU work authorization;
\item better access to European employers;
\item Portuguese language may become useful;
\item lower cost than Germany, Netherlands, or Nordics in many cities.
\end{itemize}




Weaknesses:




\begin{itemize}
\item salaries can be lower than Germany or Netherlands;
\item tax and social-security obligations reduce net income;
\item remote employment still requires payroll/EOR/contractor compliance;
\item tax incentives cannot be assumed.
\end{itemize}




\subsection{10.3 Germany resident through marriage/family route}




Best if relocation with German partner becomes real.




Strengths:




\begin{itemize}
\item direct access to German labor market after residence/work authorization;
\item high relevance to industrial visualization, simulation, digital twins, automotive, manufacturing, and XR;
\item higher salary ceiling than Portugal or Spain;
\item spouse/family route may reduce dependence on employer sponsorship;
\item stronger legal and social protections.
\end{itemize}




Weaknesses:




\begin{itemize}
\item language barrier;
\item bureaucracy;
\item health insurance and tax complexity;
\item higher cost of living;
\item family route requires correct documents and legal processing;
\item German employers may still prefer German-speaking candidates for many roles.
\end{itemize}




\subsection{11. Strategic interpretation by timeline}




\subsection{11.1 Next 0–6 months}




Primary route:




\begin{itemize}
\item Colombia-based remote contractor.
\item International applications to remote-first companies.
\item No EU work authorization claims.
\item Mention future EU mobility only when strategically useful.
\end{itemize}




Best role families:




\begin{itemize}
\item Real-Time 3D / Interactive Technical Visualization Developer;
\item Unity WebGL / Interactive 3D Developer;
\item Unity Technical Artist with optimization/runtime angle;
\item Tools / Python automation only as secondary.
\end{itemize}




\subsection{11.2 Next 6–18 months}




Primary route:




\begin{itemize}
\item prepare EU-ready portfolio;
\item track Portuguese citizenship process;
\item start building Portugal/Germany professional networks;
\item learn enough German to reduce Germany friction;
\item create EU-compatible CV and LinkedIn variants;
\item monitor Blue Card salary-fit roles.
\end{itemize}




\subsection{11.3 After Portuguese passport}




Primary route:




\begin{itemize}
\item apply as EU citizen;
\item include EU work authorization clearly;
\item target Germany, Netherlands, Ireland, Portugal, Spain, Nordics;
\item consider relocation only if net income and career signal justify it.
\end{itemize}




\subsection{11.4 If marriage/Germany route becomes active}




Primary route:




\begin{itemize}
\item verify documents with German mission and local immigration authority;
\item prepare language/administrative requirements;
\item evaluate relocation budget;
\item apply to Germany-based roles once work authorization pathway is concrete;
\item consider German-resident remote roles after compliance review.
\end{itemize}




\subsection{12. Decision thresholds}




Do not relocate to Europe only for a marginal salary improvement.




Relocation becomes rational if at least one condition is true:




\begin{itemize}
\item net income after tax and rent is materially better than Colombia;
\item role is a strong career signal;
\item role gives EU industry experience;
\item role is in digital twin, simulation, XR, industrial visualization, or technical art;
\item partner/family relocation makes Germany strategically necessary;
\item EU citizenship is already completed and relocation friction is low.
\end{itemize}




Relocation is weak if:




\begin{itemize}
\item salary is close to Colombia remote contractor income;
\item role is low-skill 3D production;
\item employer requires expensive relocation without support;
\item German bureaucracy is unresolved;
\item the role is not aligned with TwinSight / real-time 3D positioning;
\item job requires German fluency Alexander does not yet have.
\end{itemize}




\subsection{13. Open questions requiring professional verification}




Before making binding decisions, verify:




\begin{enumerate}
\item Exact status and timing of Portuguese citizenship/passport.
\item Whether Portuguese citizenship is already approved, pending, or only expected.
\item Whether Portuguese tax incentives apply to Alexander’s future profile.
\item Whether work from Portugal would be employee, contractor, or freelancer.
\item Whether the German spouse route requires A1 German in the specific case.
\item Which German mission handles the visa if applying from Colombia.
\item Whether the marriage route applies before or after the German partner has German residence established.
\item Whether a Germany-based role can start before residence title approval.
\item Whether a foreign employer can hire Alexander through EOR in Germany or Portugal.
\item Net salary simulation with a local accountant before relocation.
\end{enumerate}




\subsection{14. Legal-risk rules for the career strategy}




\subsubsection{Do not claim now}




\begin{itemize}
\item EU work authorization.
\item German work authorization.
\item Portuguese citizenship already completed.
\item ability to relocate immediately without sponsorship.
\item Blue Card eligibility without job offer and salary match.
\end{itemize}




\subsubsection{Safe to claim now}




\begin{itemize}
\item Colombia-based remote contractor availability.
\item openness to EOR if supported.
\item future EU mobility expected through Portuguese citizenship.
\item willingness to relocate under the correct legal route.
\item German partner context only if relevant and disclosed carefully.
\end{itemize}




\subsubsection{Safe to claim after passport}




\begin{itemize}
\item EU citizen.
\item authorized to work in EU/EEA.
\item available for EU employment.
\item open to relocation/hybrid roles in Europe.
\end{itemize}




\subsubsection{Safe to claim after Germany residence title}




\begin{itemize}
\item legally resident in Germany.
\item authorized to work in Germany.
\item available for German payroll or local contracts.
\end{itemize}




\subsection{15. Practical action checklist}




\subsection{Now — Colombia phase}




\begin{itemize}
\item Keep applying as Colombia-based remote contractor.
\item Do not claim EU work authorization.
\item Prepare an EU-ready CV variant.
\item Track Portugal citizenship timeline.
\item Maintain documentation for German route.
\item Start German language study only as strategic background, not as the main 3–6 month job-search lever.
\end{itemize}




\subsection{6–12 months}




\begin{itemize}
\item Prepare relocation budget scenarios.
\item Research Portugal accountant options.
\item Research Germany immigration-law consultation.
\item Build EU target-company list.
\item Start networking with Germany/Portugal technical visualization companies.
\item Create a German-market LinkedIn keyword variant.
\end{itemize}




\subsection{Once Portuguese passport is issued}




\begin{itemize}
\item Update LinkedIn work authorization.
\item Update CV header and availability line.
\item Apply to EU roles as authorized candidate.
\item Consider Portugal as administrative base.
\item Consider Germany, Netherlands, Ireland, and Nordics for higher upside.
\end{itemize}




\subsection{If marriage/Germany route activates}




\begin{itemize}
\item Confirm documents with German authorities.
\item Verify visa/residence requirements.
\item Budget for relocation and health insurance.
\item Prepare German-language basics.
\item Apply to Germany roles once timeline is credible.
\item Avoid signing employment contracts before legal start eligibility is clear.
\end{itemize}




\subsection{16. Final mobility conclusion}




The strongest mobility logic is staged.




For the next 3–6 months, the practical route is \textbf{remote international contracting from Colombia}. This preserves flexibility, avoids immigration delays, and lets Alexander monetize TwinSight while completing thesis defense and portfolio packaging.




The Portuguese passport is a high-value future asset. Once obtained, it changes the strategy materially because Alexander becomes an EU citizen and can compete for EU roles without work-permit friction.




Germany is the highest-upside European route if the relationship/marriage pathway becomes concrete. It is especially relevant for industrial visualization, digital twins, simulation, XR, and Unity-based enterprise applications. But Germany should be treated as a structured relocation pathway, not as an immediate assumption.




The correct recruiter narrative is therefore:




\begin{itemize}
\item now: Colombia-based contractor;
\item later: EU citizen once Portuguese passport is issued;
\item conditional: Germany-resident worker if family reunification through marriage is completed;
\item fallback: Blue Card only if employer sponsorship and salary thresholds align.
\end{itemize}




\subsection{17. Sources}




\begin{itemize}
\item Your Europe — Work permits: \url{https://europa.eu/youreurope/citizens/work/work-abroad/work-permits/index_en.htm}
\item Your Europe — Registering residence abroad: \url{https://europa.eu/youreurope/citizens/residence/documents-formalities/registering-residence/index_en.htm}
\item Make it in Germany — EU Blue Card: \url{https://www.make-it-in-germany.com/en/visa-residence/types/eu-blue-card}
\item Make it in Germany — Spouses joining EU citizens: \url{https://www.make-it-in-germany.com/en/living-in-germany/family-reunification/spouses-eu-citizens}
\item Make it in Germany — Visa for self-employment / freelancing: \url{https://www.make-it-in-germany.com/en/visa-residence/types/self-employment}
\item Make it in Germany — Health insurance: \url{https://www.make-it-in-germany.com/en/living-in-germany/money-insurance/health-insurance}
\item German Residence Act, section 4: \url{https://www.gesetze-im-internet.de/englisch_aufenthg/englisch_aufenthg.html}
\item PwC Germany — Individual residence: \url{https://taxsummaries.pwc.com/germany/individual/residence}
\item PwC Germany — Individual taxes on personal income: \url{https://taxsummaries.pwc.com/germany/individual/taxes-on-personal-income}
\item PwC Germany — Other taxes: \url{https://taxsummaries.pwc.com/germany/individual/other-taxes}
\item PwC Portugal — Individual residence: \url{https://taxsummaries.pwc.com/portugal/individual/residence}
\item PwC Portugal — Individual taxes on personal income: \url{https://taxsummaries.pwc.com/portugal/individual/taxes-on-personal-income}
\item PwC Portugal — Other taxes: \url{https://taxsummaries.pwc.com/portugal/individual/other-taxes}
\end{itemize}



